\begin{recomendations}
\label{ref:recomendations}

Frente a las limitaciones identificadas, las siguientes acciones permitirán llevar el sistema a su pleno potencial operativo:

\begin{enumerate}
    \item \textbf{Migrar los canales de notificación a producción.} Reemplazar EmailJS y OpenWA por un servidor SMTP institucional y la API oficial de WhatsApp Business, incorporando SMS como canal complementario para usuarios sin datos móviles.

    \item \textbf{Completar las pruebas de usabilidad con médicos y jefes municipales.} Aplicar el recorrido cognitivo y el cuestionario SUS a al menos 3 participantes de cada perfil en sesiones presenciales, registrando tiempos de ejecución y tasa de errores.

    \item \textbf{Replicar las pruebas de rendimiento en infraestructura institucional.} Ejecutar nuevamente las pruebas de carga sobre los servidores físicos del IFV para obtener umbrales reales, sin las limitaciones del entorno de desarrollo.

    \item \textbf{Optimizar los recursos estáticos del frontend.} Aplicar las mejoras detectadas por Lighthouse: compresión de imágenes a formatos WebP, eliminación de JavaScript no utilizado y precarga de estilos críticos.

    \item \textbf{Implementar fusión inteligente de duplicados.} Al confirmar un duplicado, preservar los datos únicos de la copia en el reporte original, siempre bajo supervisión del jefe de vacunación.

    \item \textbf{Migrar a servidores físicos del Instituto Finlay.} Trasladar el sistema desde Railway y Vercel a los servidores institucionales utilizando la misma configuración de contenedores Docker, garantizando el control sobre los datos de farmacovigilancia.

    \item \textbf{Desarrollar una aplicación móvil nativa.} Evaluar la creación de una aplicación para iOS y Android que aproveche notificaciones push y acceso sin conexión, apoyándose en la API actual sin requerir modificaciones.
\end{enumerate}

\end{recomendations}